\section{The class and the cup-product formula}
\label{sec:cup}

\subsection{Transition functions}

On \(W_i\cap W_j\), define

\[
z_{ij}=s_i+s_j\in\F,
\]

using any representative of the orbit. This is well defined: flipping
both bits adds \(1+1=0\). It is locally constant, as is visible on
the compact-open cells. For three indices,

\[
z_{ij}+z_{jk}=z_{ik}.
\]

Thus \(z=(z_{ij})_{i<j}\) is a degree-one cocycle in the complex
of \cref{sec:resolution}. Let

\[
\eta_r=[z]\in H^1(U_r;\F).
\]

Geometrically, \(z_{ij}\) records whether the gauge choices on the
two charts agree or differ by the involution. This observation
motivates the cocycle; the proof uses the displayed functions and
their identities directly.\leanmark{cocycle}

\subsection{Why multiplication is computed by adjacent differences}

\begin{proposition}\label{prop:cup-formula}
For \(1\leq k\leq r\), the class \(\eta_r^k\) is represented on
the face \([i_0,\ldots,i_k]\) by

\[
z_{i_0i_1}z_{i_1i_2}\cdots z_{i_{k-1}i_k}.
\]

In particular, on the full intersection and pulled back to configurations
in which every coordinate is finite, \(\eta_r^r\) is represented by

\begin{equation}\label{eq:top-cochain}
F_r((a_i,s_i)_{i=0}^r)=\prod_{j=0}^{r-1}(s_j+s_{j+1}).
\end{equation}
\end{proposition}

The familiar ordered \v{C}ech cup formula suggests this answer. We give
a chain-level proof, which also fixes its relation to ordinary sheaf
cohomology and to the operation used in the formalization.

\begin{proof}
Use the resolution \(P_\bullet^{\F}\) in sheaves of \(\F\)-modules.
Its graded Ext group

\[
\Ext^*_{\Sh(U_r;\F\text{-}\mathrm{Mod})}
(\underline\F,\underline\F)
\]

has the Yoneda product: multiplication is composition of the
corresponding morphisms in the derived category. This is the standard
realization of the cup algebra of derived global sections.

We compare this description with the ordinary cohomology already
computed. The integral and module resolutions are each independently
projective and exact. Their Hom complexes, with constant \(\F\)
as target, identify with the same complex of sections on chart
intersections. The identifications commute with restrictions and
hence with differentials. At degree zero the rank-one generator maps
to the constant function \(1\). This section-induced comparison
identifies the module Ext description with ordinary
\(H^*(U_r;\F)\), with its usual cup multiplication.\leanmark{comparison}

For a direct product calculation, write
\(d_k:P_{k+1}^{\F}\to P_k^{\F}\), and define

\[
L_k:P_{k+1}^{\F}\longrightarrow P_k^{\F}
\]

on a face \([i_0,\ldots,i_{k+1}]\) as follows: multiply by
\(z_{i_0i_1}\), delete the first vertex, and map into the summand
indexed by \([i_1,\ldots,i_{k+1}]\). A locally constant scalar
acts on sheaf sections by gluing the constant scalar actions on its
open fibers, so this construction defines a sheaf morphism.

The cocycle identity gives the chain-lift relation

\[
L_k\circ d_{k+1}=-d_k\circ L_{k+1}.
\]

Indeed, terms deleting a vertex after the first two cancel by the
face identities. The remaining terms combine by
\(z_{i_1i_2}+z_{i_0i_2}=z_{i_0i_1}\). In characteristic two the
minus sign is the same as a plus sign. If
\(\epsilon:P_0^{\F}\to\underline\F\) is the augmentation,
then \(\epsilon\circ L_0\), under the section comparison, is
precisely \(z\). Thus \(L\) lifts the degree-one cocycle to a
shifted map of the resolution.

Set \(c_0=\epsilon\) and \(c_{k+1}=c_k\circ L_k\).
The chain identity shows that these are cocycles. To see which classes
they represent, recall that the augmentation is inverted when
passing from the resolution to its resolved object in the derived
category. In composing two morphisms represented through the
resolution, the middle augmentation and its inverse cancel.
Consequently composing with the lift \(L\) computes multiplication
by \(\eta_r\). Induction shows that \(c_k\) represents
\(\eta_r^k\) for \(k\geq1\).

Finally, each application of \(L\) removes the current first vertex
and multiplies by its transition to the next one. On a face of length
\(k+1\), the composite therefore multiplies successively by
\(z_{i_0i_1},z_{i_1i_2},\ldots,z_{i_{k-1}i_k}\). This is the
claimed formula.\leanmark{cup}
\end{proof}

The product formula gives a useful simplification: although the space
has points at infinity, the top intersection has none. Its pulled-back
cochain \(F_r\) is just a function on a product of discrete sets.
The points at infinity return only when we ask whether this function
is a sum of restrictions of continuous functions on the intersections
one degree lower.
That is where the nonvanishing argument begins.
